\section{Jeres Revy}
Når man laver revy som band bliver man før eller siden konfronteret med spændingen imellem inkrusivitet og kvalitet. Som bandboss bliver det særligt synligt hvis man deltager i materialeudvælgelse, rollefordeling og øveplanlægning.

I min optik er studenterrevyer et grundlæggende socialt projekt: De er til for at skabe sammenhold og kutur i et studiemiljø.
Derfor giver det ikke mening for en revy at være eksklusiv.
Alligevel har de fleste bands nok oplevet frustrationen ved at gøre sig umage for at lære et nummer, både individuelt og som band, for at blive mødt af en "sanger" der hverken kan sin tekst eller melodiens toner og takt.

Det sker af og til at meget dygtige revyster, som er store ressourcer for resten af revyen, holder op med lave revy, fordi de føler at deres tid og indsats bliver spildt.
At en revy er inklusiv betyder ikke blot at man kan være med hvis man vil; det betyder også at man vil være med hvis man kan.

Det er op til hver revy at beslutte hvordan man håndterer dette, og i min tid har det også være op til hver enkelt person.
Mit indtryk er, at de fleste bare lærer at leve med det.
Det vil nok altid kræve en eller anden form for generøsitet at være band til en studenterrevy - det bliver i hvert fald vanskeligt at få scenegængere til at øve lige så meget som revyens musikere.

Min eneste konkrete idé til en forbedring (som jeg ikke selv har afprøvet) kommer fra matematikrevyens koreografer, der inddeler dansene i lette, mellem og svære.
I tilmeldingen/rolleønskningsskemaet kan dansende scenegængere så krydse af, hvilken sværhedsgrad af dans de har mod på.
Det er min forståelse at ønskning af svære danse også forstås som en accept af højere forventninger under øvning.
Måske kunne man gøre noget tilsvarende for sangrolle-ønsker?

Under alle omstændigheder kan et totalt fravalg af krav være en bjørnetjenste for de svage/uerfarne revyster.
Hvis det i sidste ende betyder, at de dygtigste revyster forlader revyen, bliver det meget mindre sjovt at lave revy for de svage revyster.
Det er sjovt at være på scenen til et nummer der er godt.
Det er mindre sjovt at være på scenen til et nummer der er dårligt.

En ideel model ville både have plads til at alle kan deltage og bidrage i det omfang de magter, og stadig levere tilfredsstillende kreative oplevelser for de dygtige og dedikerede.

Jeg synes man burde overveje, og snakke om (i rette forum):
\begin{itemize}
    \item Hvorfor laver vi revy?
    \item Hvor vigtigt er det sociale projekt?
    \item Hvor vigtigt er det at lave et godt show?
    \item I hvilken grad kan vi stille krav til hinanden?
\end{itemize}